\documentclass[10pt]{article}
\usepackage[a4paper,margin=23mm]{geometry}
\usepackage{amsmath,amssymb,booktabs,graphicx,hyperref}
\setlength{\emergencystretch}{2em}
\title{From Partial Evidence to Complete 3D Responses:\\Query-Conditioned Evidence Fusion}
\author{TMM working manuscript --- controlled validation in progress}
\date{October 18, 2026}
\begin{document}
\maketitle
\begin{abstract}
Visual teachers provide local evidence for scene queries, but a strong match
does not establish an object's complete extent or authorize exclusion of its
surroundings. We study how such evidence can form a complete three-dimensional
response without conflating missing information with background. We formulate
query-conditioned evidence fusion: determine which observations refer to the
target, assign their support and exclusion scope, and combine admitted responses
through geometry with an output-appropriate objective. An ambiguity argument,
an exact consensus-attenuation result and an occlusion counterexample explain
why evidence scope and output semantics matter even with fixed features and
geometry. The method retains native teacher matching and provides typed
semantic and complete-membership realizations. For reference queries, it
completes and admits the reference response before cross-view transport.
Evaluation covers five task outputs and 31 scene jobs. Matched component
controls support identity admission, response integration and complete extent;
in an eight-task NVOS source-only control, reference completion improves IoU
from 88.01\% to 92.75\% under the same downstream fusion. A separately declared
target-image-conditioned extension reaches 92.94\% with its primary feasibility
policy and 93.17\% IoU with 98.95\% AP using its fixed source-quality secondary
policy. Lower source loss alone does not consistently improve task accuracy.
Input differences and the opened development cohort delimit the empirical
claims; the analysis does not guarantee correctness of estimated evidence scope.
\end{abstract}
\section{Introduction}
A scene query specifies both an identity and the spatial extent to be returned.
A text description may match several observations, a reference image may expose
only part of an object, and a point identifies one surface rather than every
member. Modern visual teachers provide strong local correspondences. Turning
these correspondences into a complete three-dimensional answer still requires
deciding what each observation actually establishes.

A reliable match and a complete observation are different properties. A mask
may correctly identify a flower while omitting its leaves. Transporting that
mask accurately can reproduce the omission across views; treating its exterior
as background can actively suppress the missing members. Conversely, expanding
the response without preserving identity can capture a neighboring object.
The scientific question is therefore: \emph{when does partial query evidence
justify support or exclusion, and how should this evidence form a complete
response under the intended output semantics?}

We address this question with \emph{query-conditioned evidence fusion} (QCEF).
The central mechanism assigns evidence authority before aggregation: identity
applicability determines which observations can act, and spatial scope
determines where they can support, exclude or leave the response unknown.
Native teacher scores remain in their own matching spaces. Admitted responses
then enter a geometric fusion or conditional-membership operator appropriate
to the requested output. This is a shared evidence-to-response formulation
with typed realizations, not a claim that every modality uses one network or
one inverse solver.

Reference queries make the mechanism particularly concrete. We complete the
reference response using signed query evidence, select between whole reference
candidates using reference-side feasibility or predicted quality, and only then
transport the selected response to other views. This changes the scope of the
transmitted evidence rather than asking a stronger reconstruction objective to
repair incomplete targets. For map-only image and point queries, the analogous
extent requirement is implemented by conditional membership over the complete
carrier. Complete here describes the output domain and intended capability;
it does not imply that every predicted object is recovered without error.

The analysis distinguishes three obstacles. Partial labels alone need not
identify complete extent. Incorrect negative authority attenuates a weighted
response even when correspondences are correct. Finally, rendering a selected
set and averaging membership in the full scene can disagree under occlusion.
These statements motivate scope estimation and typed objectives, without
turning fallible teacher predictions into guaranteed semantic truth.

Our implementation includes the established system recorded as v10 and the
reference-completion extension. Its evidence spans LERF semantic segmentation
and projected selection, ScanNet open-vocabulary segmentation, NVOS reference
segmentation and multimodal UQIS. Historical controls isolate identity,
response-use and membership capabilities. Matched reference interventions show
that completing scope before transport is useful even without target-image
evidence; stronger source fitting alone has mixed effects. Published results
and verified local reproductions provide comparison context with their input
and metric differences retained.

Our contributions are:
\begin{enumerate}
\item A response-level formulation of partial-evidence authority, with scoped
analysis of extent ambiguity, negative-evidence attenuation and output
non-equivalence. It separates the information required to identify an object
from the evidence needed to recover its members.
\item A query-conditioned fusion method that retains native matching while
combining admitted evidence through typed geometric and membership operators.
Reference completion and whole-response admission before transport provide a
concrete mechanism for repairing incomplete query scope.
\item Five-output evaluation with matched component interventions, protocol
audits and counterexamples to stronger-fitting explanations. The evidence
connects measured gains to specific capabilities and records where the shared
formulation has not yet yielded a universally effective admission rule.
\end{enumerate}

\section{Related work}
\paragraph{Lifting visual evidence to scene representations.}
VALA~\cite{vala} uses visibility-aware feature lifting and robust multi-view
aggregation for open-vocabulary Gaussian segmentation. LUDVIG~\cite{ludvig}
aggregates lifted visual features and applies graph diffusion; Occam's
LGS~\cite{occam} provides an efficient Gaussian language representation.
These approaches establish the value of correspondences, visibility and
aggregation. We keep the available teacher representations and focus on
forming query responses from them. Neither weighted averaging, robust fitting
nor frozen native features alone is claimed as our contribution.

\paragraph{Query-dependent evidence and complete extent.}
EviSplat~\cite{evisplat} preserves observation features in class-agnostic
instances and combines query-selected evidence with local Gaussian support.
Query-dependent use of preserved evidence is therefore prior art. Our focus is
the distinction between identity applicability and authority to support or
exclude spatial extent, including completion of reference scope before
transport and compatibility of response objectives with actual selection.
The distinction is evaluated through matched interventions; the formulation
does not establish that our scope estimates are universally more accurate.

\paragraph{Teacher capabilities and response-level adaptation.}
SAM3~\cite{sam3}, SigLIP2~\cite{siglip2}, DINOv2~\cite{dinov2} and
Mask-Adapter~\cite{maskadapter} supply attributed segmentation or matching
capabilities. We study how their evidence is used, not a new foundation model.
For reference queries, frozen image decoding and tracking are combined with
reference-side admission and geometric response fitting. For map-only queries,
complete membership includes independent labelled-room supervision. Published
method comparisons and retained local reproductions are contextual when
teacher access, prompts, geometry or evaluation domains differ; our component
controls provide the narrower causal evidence.

\section{Problem formulation and analysis}
\label{sec:analysis}
Let $\mathcal C=\{1,\ldots,N\}$ be a fixed scene carrier and $q$ a query.
The ideal response $r_q^*$ specifies target membership or semantic relevance
over that carrier. An observation supplies a bounded response $y_e$, a geometric
observation operator $A_e$, and query-dependent information about its scope.
A deployed answer is $\widehat a_q=\mathcal R_t(r_q)$, where the output type $t$
determines the readout $\mathcal R_t$. The evidence-to-response problem is to
estimate $r_q$ from the legally available observations without silently adding
negative information or changing the meaning of $\mathcal R_t$.

We distinguish \emph{applicability}, whether an observation refers to the
requested identity, from \emph{spatial authority}, the locations at which it
can provide support or exclusion. Reliability weights the strength of already
admitted evidence. These quantities need not be calibrated probabilities.
Unknown locations supply neither positive nor negative labels; an output
adapter may still need a declared default where no observation is available.

\paragraph{Proposition 1: partial scope does not identify complete extent.}
Suppose the only constraints on binary membership are $r_i^*=1$ for $i\in P$
and $r_i^*=0$ for $i\in N$, with disjoint $P,N$ and a nonempty unknown set
$U=\mathcal C\setminus(P\cup N)$. If no other observation or constraint
distinguishes assignments on $U$, complete membership is not identifiable.

\emph{Proof.} Pick $j\in U$. Two assignments can agree everywhere except at
$j$, taking values zero and one there, while satisfying every supplied
constraint. An estimator given identical evidence returns the same answer
for both and cannot be correct on both. $\square$

This elementary ambiguity result does not assert that the full RGB input is
uninformative. Additional image evidence, native correspondence or learned
membership can resolve an ambiguity. It shows why the complement of a partial
positive observation is not, by itself, an observed negative. Our estimated
completion supplies a hypothesis about that missing scope, not a proof of it.

\paragraph{Proposition 2: false exclusion attenuates direct consensus.}
For one carrier element, let the effective evidence and fixed-prior mass be
$D>0$, numerator $S$, and response $r=S/D\in[0,1]$. Add incorrectly authorized
negative evidence with response zero and total effective mass $u\geq0$,
holding all previous weights, priors and correspondences fixed. Then
\begin{equation}
 r'=\frac{S}{D+u},\qquad r-r'=\frac{ru}{D+u}.
 \label{eq:attenuation}
\end{equation}
If a positive decision uses $r\geq\tau>0$, it flips to negative exactly when
$u>D(r/\tau-1)$.

\emph{Proof.} Zero-valued evidence adds $u$ to the denominator and nothing to
the numerator. Subtraction gives Equation~\ref{eq:attenuation}; rearranging
$S/(D+u)<\tau$ gives the threshold condition. $\square$

The result applies to the direct consensus operator with fixed existing mass,
not to every robust solver or learned member decoder. It holds with correct
geometry: better alignment cannot repair invalid exclusion authority. It also
does not justify deleting valid negatives. Signed-hint controls test that
separate issue empirically.

\paragraph{Counterexample: membership and semantic rendering do not commute.}
Consider a ray with an unselected front Gaussian of opacity $a\in(0,1)$ and
a selected rear Gaussian of opacity $b\in(0,1)$. The normalized full-scene
member response and the opacity of the selected set are respectively
\begin{equation}
 s_{\mathrm{full}}=\frac{(1-a)b}{a+(1-a)b},\qquad
 \alpha_{\mathrm{selected}}=b.
 \label{eq:readout-counterexample}
\end{equation}
With $a=0.9,b=0.2$, they are approximately $0.0217$ and $0.2$.
They produce different answers under the retained semantic threshold $0.5$
and selected-opacity threshold $10/255$. Removing the front Gaussian changes
transmittance; it is not just thresholding a value in the original renderer.
Thus preserving a selection capability requires preserving its objective and
readout, not only its underlying features. This is a constructive counterexample,
not an accuracy guarantee or a claim that all rays disagree.

\paragraph{Consequences for method and evaluation.}
These statements motivate three linked decisions within one mechanism:
admit evidence for the requested identity, estimate its support/exclusion
scope, and form the response with a compatible forward model and readout.
They predict that reconstructing incorrect evidence more accurately can retain
or amplify its errors. They do not predict the sign or magnitude of a benchmark
gain. We test the corresponding design choices with fixed-evidence solver
controls, matched reference-scope interventions and separate semantic/selection
evaluation, rather than treating lower source loss as evidence of correctness.

\section{Query-conditioned evidence fusion}
\label{sec:method}
The method forms a response in three stages: determine which observations
apply to the query, assign their spatial support and exclusion authority, and
combine their admitted responses through geometry. The reference extension repairs scope before transport. Semantic and
complete-instance outputs retain their own objectives and readouts. The following equations
describe the implemented system and identify its query-specific realizations.

\begin{figure}[ht]
\centering\includegraphics[width=\linewidth]{figures/method_overview.pdf}
\caption{Our evidence-to-response method. Applicability and spatial authority
determine how native evidence acts on the response. Geometric aggregation and
complete membership have distinct typed implementations. Reference completion precedes cross-view transport. Actual selections are
evaluated separately from continuous semantic responses.}
\label{fig:method}
\end{figure}

\subsection{Native evidence and query applicability}
For query $q$, observations $\mathcal E_q$ supply responses $y_e\in[0,1]$,
nonnegative geometric operators $A_e$ and visible domains $v_e$. Nonempty
operator rows are normalized. Text similarities and local visual correspondence
stay in their native teacher spaces; points and reference masks retain their own
semantics. We combine response evidence, not raw feature vectors. LERF, NVOS
and ScanNet use declared source-assisted observations; UQIS reads its map without
source RGB or depth at query time.

Let $g_e(q)$ express whether observation $e$ applies to the requested target.
For text candidates, we compare one- and two-component native-score models
within candidates already detected for the same query. If their Bayesian
information criteria are $B_1,B_2$, the implemented support is
\begin{equation}
 a_e=(1-\pi_2)+\pi_2\gamma_{\rm high}(s_e),\quad
 \pi_2=\sigma((B_1-B_2)/2),\quad g_e=\mathbf1[a_e\geq0.5],
 \label{eq:admission}
\end{equation}
where $\gamma_{\rm high}$ is responsibility for the higher-score component.
A single mode supplies no relative identity veto. Degenerate or missing evidence
does not invent one. Rejected regions become unknown rather than background.
Mixture support is not a calibrated probability of semantic truth.

This rule has a controlled LERF result. Reference queries instead inherit
identity through a positive-prompt mask and source tracks; categories combine
native probabilities with literal-category evidence. Image queries use local
correspondence to select identity hypotheses. Actual points supply explicit
spatial conditions. These are typed estimators, not unreported copies of
Equation~\ref{eq:admission}.

\subsection{Spatial authority: support, exclusion and unknown}
An applicable observation need not describe complete extent. Let $p_e,n_e$
denote disjoint support and exclusion domains and $c_e^+,c_e^-$ their weights:
\begin{equation}
 w_e=v_e g_e(q)(c_e^+p_e+c_e^-n_e).
 \label{eq:authority}
\end{equation}
Unobserved space and rejected detections carry no exclusion solely because they
lack a match. Predicted reference-mask complements and category competition can supply
exclusion under their stated completeness assumptions. These are fallible
provider estimates, not hard negative user prompts. Our LERF path vetoes negatives where
native positive support exists. NVOS uses interior positive/negative regions and
tracked negative reliability. ScanNet keeps both native priors and source weights.

\subsection{Reference-scope completion and admission before transport}\label{sec:reference-completion}
The NVOS extension estimates a complete reference response before cross-view
transport. For reference image $I_r$, let $\mathcal P,\mathcal N$ be original
positive and negative scribble patches. A patch containing both signs is unknown;
unmarked patches do not supply negative labels. With frozen normalized DINOv2
descriptors $f_r$, define
\[
 h_r(u)=\max_{v\in\mathcal P} f_r(u)^\top f_r(v)
       -\max_{v\in\mathcal N} f_r(u)^\top f_r(v),
 \qquad a_r(u)=\sigma(h_r(u)/\tau_r).
\]
The source-only scale $\tau_r$ is the median signed leave-one-out reference
margin divided by $\log9$. Missing signed support or a nonpositive scale falls
back to the original reference mask. Upsample $a_r$ and decode a semantic reference response with frozen SAM3,
using its thresholded extent as a box and its log-odds as a dense prompt.
Probabilities are clipped to $[10^{-4},1-10^{-4}]$ before log-odds conversion;
the dense prompt uses the decoder's $288\times288$ grid. A second candidate additionally uses 16 deterministically spread points
per original scribble sign. Threshold both at 0.5; for the constrained candidate,
impose all original positive/negative pixels exactly, without relabelling
unknown pixels. Agreement with sparse annotations does not prove complete scope.

Two fixed reference-only admission policies are compared. The feasibility policy
retains the semantic candidate when it satisfies every original signed pixel,
otherwise choosing the constrained candidate. The source-quality policy chooses
\[
 \hat k=\arg\max_{k\in\{\mathrm{semantic},\mathrm{constrained}\}}
             \widehat Q_r(k),
\]
where $\widehat Q_r$ is SAM3's predicted mask quality at reference decoding,
before the final hard annotation constraints. Ties prefer the semantic candidate;
missing quality falls back to feasibility. These scores are fallible predictions,
not target GT or calibrated correctness probabilities. Neither rule reads a
scene identifier, target score or task-specific threshold.

The selected whole reference mask initializes frozen two-frame SAM3 memory,
producing $y_t=T(I_r,m_r^{\hat k},I_t)$. The native-ray fit below then
creates and actually re-renders the 3D response. Original negative scribbles,
DINOv2 and observed target RGB remain declared extensions to v10's input contract.
This changes response construction and admission, keeping teacher weights and
Gaussian geometry fixed. Both fixed policies are reported on all eight tasks;
reference-admission design follows the opened-cohort completion diagnosis.
A source-only control instead propagates original or completed references
through the same source-view sequence and fits a response without target-image
evidence. This tests whether the scope intervention survives multi-view fusion.

\subsection{Geometric response formation}
For direct consensus, define $m=\sum_e A_e^\top w_e$ and
$b=\sum_e A_e^\top(w_ey_e)$. The response is
\begin{equation}
 r_i=\begin{cases}
 (b_i+\lambda_i\bar r_i)/(m_i+\lambda_i),&m_i>10^{-8},\\
 r_{0,i},&\text{otherwise}.
 \end{cases}
 \label{eq:consensus}
\end{equation}
LERF uses $\lambda=0$. ScanNet uses a two-native mean $\bar r$ and prior strength
0.1 times visible mass. This is an established weighted-consensus operation;
the tested design concerns admitted evidence and its actual use in the response.
The active ScanNet implementation has no fitted local residual head.

The original NVOS path instead minimizes robust rendered discrepancy,
\begin{equation}
 \min_{r\in[0,1]^N}\sum_e\langle w_e,\rho_{0.25}(A_er-y_e)\rangle,
 \label{eq:reference}
\end{equation}
with 120 mass-preconditioned updates from a neutral 0.5 response. Its output
adapter zeros unsupported elements. Changing this solver while preserving
evidence is an ablation. Our controls do not support replacing every typed
operator by one inverse solver.

\paragraph{Native-ray fitting of completed reference responses.}
Let $p_t$ denote the transported target-image response and $s^0$ the original
v10 Gaussian response. Native-ray writeback retains
$y_t=p_t$ on the
$1008\times756$ output grid and uses $w_t=|2p_t-1|v_t$. Here $v_t$ admits rays
with nonzero observed mass and more than half of their rendered alpha supplied
by the original observed Gaussians. The normalized observation matrix
$\bar A_t$ is formed on the same ray grid, using the fixed original Gaussian
geometry. Initialize at $s^0$ and take 120 bounded Huber updates for
\begin{equation}
 \min_{0\leq s\leq1}\sum_u w_{tu}\rho_{0.25}((\bar A_t s)_u-y_{tu}),
 \qquad s_i=s_i^0\ \text{when no admitted observation supports }i.
\end{equation}
We measure both $p_t$ and the actual re-rendered 3D response. Geometry and
teacher weights remain fixed. Because this update uses view $t$'s image,
the result is view-conditioned; a query-global field and preserved source-view
accuracy do not follow from target-view gains. Original source-teacher
residuals are reported as a separate compatibility check.

\subsection{Complete extent and typed outputs}
For source-assisted actual selections, our implementation also fits membership
against the silhouette of the selected Gaussians. This is distinct from rendering
their semantic scores in the full scene. Let $z_i\in[0,1]$ be membership and
$o_{vi}(u)$ the uncomposited per-Gaussian opacity at pixel $u$. The selected
silhouette has the compositing form
\begin{equation}
 \alpha_v(u;z)=1-\prod_i(1-z_i o_{vi}(u)).
\end{equation}
The implemented loss combines foreground/background-balanced binary
cross-entropy and soft Dice on admitted source pixels, plus a 0.1-weighted
projected-footprint membership penalty. The footprint is normalized by its
cost at source-initialized membership. Optimization uses at least 240 Adam
updates (extended to the source-view count when larger) at learning rate 0.05,
cycling through the active views. Actual membership is then
thresholded at 0.5 and its selected opacity is rendered. The existing projection
readout thresholds quantized 8-bit alpha above 10. This describes an existing
v10 capability, not a new solver introduced by the subsequent NVOS experiment.

\paragraph{Native image identity and conditional complete membership.}
For an image query, let $q_u,f_v$ be normalized query and regional patch
features, with unit-sum weights $a_u,b_v$. Local correspondence uses the
symmetric weighted Chamfer score
\begin{equation}
 k(Q,R)=\tfrac12\left[\sum_{u\in Q}a_u\max_{v\in R}q_u^\top f_v+
                         \sum_{v\in R}b_v\max_{u\in Q}q_u^\top f_v\right].
\end{equation}
The deployed image route averages the candidate ranks from native SigLIP
matching and local DINOv2 correspondence. It selects a nonempty region, then
uses its actual member nearest the coordinate-wise 3D median as the identity
anchor. Both teacher spaces remain unchanged. An actual point supplies a
geometric anchor directly; an image-derived anchor remains an identity estimate.

For anchor $a$ and every carrier element $i$, the complete-membership decoder
predicts
\begin{equation}
 p_i=f_\theta(\phi(i,a)),\qquad S_a=\{i:p_i>\tau\},\quad p_a=1.
\end{equation}
The 15-dimensional evidence vector retains view-balanced nested-mask votes,
any-mask/all-mask agreement and their disagreement, common visibility and its
presence flag, anchor/element visibility fractions, bounded signed displacement
and distances, native appearance agreement, and native-support availability.
Nested masks share a view's vote instead of becoming independent observations.
The decision is evaluated over the complete carrier, without requiring an
intermediate mask to be a pure or complete instance.

The deployed decoder is an Extra Trees classifier with 64 trees, at most 512
leaves per tree and minimum leaf size 32. It is inherited from training on
independent 3D instance labels in 22 real-observation rooms; three separate
calibration rooms select $\tau=0.35$. Neither split overlaps the benchmark
rooms. This supervised decoder is disclosed as part of the method; freezing
teacher features does not make the complete system training-free. The same
anchor-positive rule is applied to inferred image anchors. Its identity
condition is fallible, and an ablation removing only that assignment leaves
all measured selected sets unchanged.

Continuous semantic response, category argmax and actual instance selection
have different readouts. LERF and NVOS maintain independent Gaussian sets;
changing the semantic response does not change that set automatically.
Selected-set projection IoU uses 2D GT, not volumetric GT. UQIS complete
membership is not Equation~\ref{eq:consensus}. These distinctions are essential
to evaluating functional improvement.

The counterexample in Section~\ref{sec:analysis} explains why these output
operators cannot be exchanged silently. Field-by-readout controls evaluate the
actual retained NVOS fields separately; they are evidence about the output
contract, not a new gain in either field.

\section{Experiments}
\label{sec:experiments}
We organize evaluation around three questions: whether query identity must be
resolved before fusion, whether repairing reference scope improves complete
extent, and whether better reconstruction or a different readout alone explains
the gains. Comparisons establish task performance; matched interventions address
these mechanism questions.

\subsection{Protocol, inputs and supervision}
The established method covers 31 scene jobs and five task outputs. Four LERF
scenes supply two outputs, so these are not five independent datasets.
All values are percentages. The benchmark rooms have been used during
development; neither prediction-before-scoring seals nor scene-level intervals
make this an untouched final test.
\begin{table}[ht]
\centering\small
\begin{tabular}{lp{.25\linewidth}p{.43\linewidth}}
\toprule
Task & Cohort and primary metric & Input and evaluation contract \\
\midrule
LERF-2D & 4 scenes, 208 observations; mIoU and localization & Source-assisted text; semantic response rendered before thresholding. \\
LERF-3D selection & Same scenes; projection mIoU and Acc.@0.25 & Independent selected Gaussians projected to 2D GT, not volumetric GT IoU. \\
ScanNet OVS & Paper8 Gaussian mIoU/mAcc; internal10 mesh mIoU & Fixed 19/15/10 menus; Gaussian and mesh domains kept separate. \\
NVOS & 8 tasks, 7 scene clusters; IoU and AP & Original reference/source access and target-image extensions explicitly separated. \\
UQIS & 9 rooms, 67 instances, 268 queries; set IoU and AP & Text, image, 2D point and 3D point; map-only query-time access. \\
\bottomrule
\end{tabular}
\caption{Evaluation units. LERF and UQIS average within scenes then across
scenes; NVOS averages tasks. Its LUDVIG reproduction first averages three seeds
within each task.}
\label{tab:protocols}
\end{table}

The implementation recorded as v10 is our method, including native matching,
admission, fusion and complete membership. LERF uses SAM3 source masks and
SigLIP2 matching; ScanNet retains SigLIP2 and Mask-Adapter priors with
literal-category SAM3 evidence. Original NVOS uses positive-derived SAM3
reference/source tracks. All UQIS maps retain SigLIP2 giant and DINOv2 ViT-B/14.
Its complete-member decoder uses 22 independently labelled training rooms and
three calibration rooms, disjoint from the benchmark. Frozen teachers do not
make this supervised component training-free. The reference-completion
extension adds original negative scribbles and DINOv2; its target-conditioned
variant also uses observed target RGB. It does not replace the input contract
of the original five-output results.

Table origins are \textbf{P}: published, \textbf{R}: locally reproduced,
\textbf{O}: our established implementation, and \textbf{M}: a measured extension.
Published values retain their reporting-paper provenance. Local aggregates
are recomputed from hash-bound reports; no new external training is performed.
No per-scene or per-task winners are combined. The full experiment catalog,
secondary variants, both NVOS GT conventions and optional complete-response
replacement are retained in the supplements.

\subsection{Five-output performance and protocol checks}
\paragraph{LERF semantic response and actual selection.}
Tables~\ref{tab:lerf2d}--\ref{tab:lerf3d} report semantic IoU/localization and
selected-set projection separately. Our method achieves 70.40 semantic IoU
and 66.96 projection IoU. SAM3 text-conditioned evidence and native SigLIP2
matching differ from the external methods' teacher access; these differences
cannot be assigned to the admission mechanism alone. The VALA reproduction
uses compatible Occam RGB geometry rather than retraining VALA geometry.
An independent CPU audit exactly reproduces all 208 selected-mask IoUs;
semantic records are reaggregated without claiming a fresh render.
\begin{table}[ht]
\centering\small\begin{tabular}{llrr}
\toprule
Method & Origin & mIoU & Loc. acc. \\
\midrule
LERF \cite{vala} & P & 37.40 & 73.60 \\
LangSplat \cite{vala} & P & 51.40 & 84.30 \\
GAGS \cite{vala} & P & 54.10 & 81.70 \\
LangSplatV2 \cite{vala} & P & 59.90 & 84.10 \\
Occam's LGS \cite{vala} & P & 61.30 & 82.50 \\
VALA \cite{vala} & P & 61.70 & 86.40 \\
Occam's LGS & R & 63.62 & 82.85 \\
Ours (v10) & O & 70.40 & 92.42 \\
\bottomrule
\end{tabular}

\caption{LERF-2D. Published context from VALA Table 1~\cite{vala}.
Localization is peak-in-box accuracy, not ranking AP. Teacher access differs.}
\label{tab:lerf2d}
\end{table}
\begin{table}[ht]
\centering\small\begin{tabular}{llrr}
\toprule
Method & Origin & mIoU & Acc.@0.25 \\
\midrule
OpenGaussian \cite{vala} & P & 38.36 & 51.43 \\
Dr. Splat \cite{vala} & P & 43.29 & 64.30 \\
InstanceGaussian \cite{vala} & P & 43.87 & 61.09 \\
CAGS \cite{vala} & P & 50.79 & 69.62 \\
VoteSplat \cite{vala} & P & 50.10 & 67.38 \\
Occam's LGS \cite{vala} & P & 47.22 & 74.84 \\
VALA \cite{vala} & P & 58.02 & 82.85 \\
VALA (compatible geometry) & R & 54.12 & 79.35 \\
Ours (v10) & O & 66.96 & 90.02 \\
\bottomrule
\end{tabular}

\caption{LERF selected-Gaussian projection, using 2D masks as GT.
Acc.@0.25 counts projections with IoU above 0.25. Published values from
VALA Table 1~\cite{vala}; its local reproduction uses compatible RGB geometry.}
\label{tab:lerf3d}
\end{table}

\paragraph{ScanNet open-vocabulary segmentation.}
Table~\ref{tab:scannet} separates eight-scene Gaussian scoring from ten-scene
mesh scoring. Our unchanged labels achieve 66.95 mIoU19 under the paper-eight
Gaussian metric and 61.01 under ten-room mesh scoring. The former applies the
VALA evaluator with anisotropic Mahalanobis pseudo labels, five-scale radius,
top-1000 class-balanced votes and opacity--volume weights. Predictions are
sealed before target-label scoring. Independent scoring reproduces all ten
mesh scenes and eight Gaussian scenes; all 1,134,207 pseudo labels agree,
including 256 checks with the released VALA implementation. On the \emph{same}
eight rooms, mesh and Gaussian mIoU19 are 61.0567 and 66.9480: that 5.8913-point
change is caused by the metric, not a learned improvement. The metric bridge
still uses provided unlabelled mesh geometry, whereas VALA uses RGB-optimized
Gaussians. Geometry and teacher access remain comparison qualifications.
\begin{table}[ht]
\centering\small\resizebox{\linewidth}{!}{\begin{tabular}{llrrrrrr}
\toprule
Method & Origin & IoU19 & Acc19 & IoU15 & Acc15 & IoU10 & Acc10 \\
\midrule
\midrule\multicolumn{8}{l}{\textit{Eight scenes: Gaussian scoring (input differences disclosed)}} \\
LangSplat \cite{vala} & P & 2.45 & 8.59 & 3.45 & 13.21 & 6.48 & 21.89 \\
OpenGaussian \cite{vala} & P & 27.73 & 42.01 & 29.67 & 46.15 & 39.93 & 57.34 \\
Dr. Splat \cite{vala} & P & 29.31 & 47.68 & 33.25 & 54.33 & 44.19 & 65.19 \\
Occam's LGS \cite{vala} & P & 31.93 & 48.93 & 34.25 & 53.71 & 45.16 & 64.39 \\
VALA \cite{vala} & P & 32.11 & 50.05 & 35.10 & 54.77 & 46.21 & 65.61 \\
VALA & R & 34.53 & 51.59 & 37.96 & 56.77 & 47.36 & 67.47 \\
Ours (v10, metric bridge) & O & 66.95 & 77.94 & 67.86 & 79.61 & 78.07 & 87.89 \\
\midrule\multicolumn{8}{l}{\textit{Ten scenes: mesh scoring (separate evaluation track)}} \\
Ours (v10) (10-mesh) & O & 61.01 & -- & 62.53 & -- & 72.49 & -- \\
Ours + replacement (L11) (10-mesh) & M & 61.01 & -- & 62.53 & -- & 72.49 & -- \\
\bottomrule
\end{tabular}
}
\caption{ScanNet OVS. Published context from VALA Table 2~\cite{vala}.
Paper8 Gaussian and internal10 mesh scoring are separate tracks.
The metric bridge aligns scoring, not teacher or geometry access.}
\label{tab:scannet}
\end{table}

\paragraph{NVOS reference responses.}
Table~\ref{tab:nvos} shows the established response, both fixed reference
admission policies and the highest-IoU refinement as an explicit tradeoff.
The primary feasibility policy reaches 92.94 IoU; the source-quality secondary
reaches 93.17 IoU and 98.95 AP. Direct image evidence reaches 93.50 IoU but is
not substituted for the actual re-rendered 3D result. Original v10 uses cached
source evidence without online target RGB; these extensions use target RGB,
negative reference scribbles and DINOv2. The improvement over 86.51 therefore
is not an input-identical ablation.

The local LUDVIG-SAM reproduction reaches 91.26 over eight tasks and three
seeds. Its GT mask-maximum convention differs from our retained $g>127$
convention; changing only this convention takes original v10 from 86.5134 to
86.6835, far below the observed extension gain. Changing bilinear
implementations moves the result by less than 0.00002 points. These checks
limit scorer explanations but do not make different input protocols identical.
Recent published rows are contextual: VCAR uses seven targets, and LEGO's
seven-scene target roster has not been verified against ours. We do not claim
same-protocol SOTA.
\begin{table}[ht]
\centering\small\begin{tabular}{llrr}
\toprule
Method & Origin & IoU & AP \\
\midrule
NVOS \cite{ludvig} & P & 70.10 & -- \\
SA3D-TRF \cite{ludvig} & P & 90.30 & -- \\
SA3D-GS \cite{ludvig} & P & 92.20 & -- \\
SAGA \cite{ludvig} & P & 92.60 & -- \\
OmniSeg3D \cite{ludvig} & P & 91.70 & -- \\
LUDVIG-DINOv2 \cite{ludvig} & P & 92.40 & -- \\
LUDVIG-SAM \cite{ludvig} & P & 91.30 & -- \\
LUDVIG-SAM2 \cite{ludvig} & P & 91.30 & -- \\
VCAR \cite{vcar} & P & 93.50 & -- \\
LEGO \cite{lego} & P & 94.20 & -- \\
WildSeg3D \cite{wildseg3d} & P & 94.10 & -- \\
Trace3D \cite{trace3d} & P & 92.50 & -- \\
LUDVIG-SAM & R & 91.26 & -- \\
Ours (v10) & O & 86.51 & 98.02 \\
Ours + reference admission (primary) & M & 92.94 & 98.86 \\
Ours + reference admission (secondary) & M & 93.17 & 98.95 \\
Ours + local decoding (tradeoff) & M & 93.28 & 98.47 \\
\bottomrule
\end{tabular}

\caption{NVOS comparison context. P values retain LUDVIG Tables 1/D and
later papers' provenance~\cite{ludvig,vcar,lego,wildseg3d,trace3d}.
N14 policies are fixed across all eight tasks. N16 has the highest measured
native IoU but lower AP (98.47 versus 98.95); it is not a uniform improvement.
All intervening variants remain in the diagnostic supplement.}
\label{tab:nvos}
\end{table}

\paragraph{Multimodal UQIS.}
Table~\ref{tab:uqis} evaluates the local ScanNet-UQIS9 v0.3.1 protocol. None
of the cited external tables evaluates these same 268 queries, so we do not
fabricate an external row. Regional readout provides an internal capability
reference; the paired component experiments below provide finer attribution.
An optional complete-response selector improves image IoU by 1.088 points,
with one query improved and other outputs unchanged. Its controls and limited
transfer are documented separately, rather than used to explain v10's gains.
\begin{table}[ht]
\centering\small\resizebox{\linewidth}{!}{\begin{tabular}{llrrrrrrrr}
\toprule
Method & Origin & T-IoU & T-AP & I-IoU & I-AP & P2-IoU & P2-AP & P3-IoU & P3-AP \\
\midrule
Ours (regional readout) & O & 12.11 & 23.69 & 26.55 & 37.75 & 42.09 & 52.74 & 41.52 & 52.36 \\
Ours (v10) & O & 16.39 & 26.60 & 39.90 & 59.31 & 49.74 & 75.79 & 49.05 & 75.33 \\
\bottomrule
\end{tabular}
}
\caption{All four UQIS modalities: text (T), image (I), 2D point (P2) and
3D point (P3). Actual-set IoU and continuous-response AP are distinct.
Both rows use our system; the regional baseline is not an external method.}
\label{tab:uqis}
\end{table}

\subsection{Q1: does query identity need to precede response formation?}
\begin{table}[ht]
\centering\small\resizebox{\linewidth}{!}{\begin{tabular}{lrrr}
\toprule
Controlled intervention & Before & After & $\Delta$ (pp) \\
\midrule
Query-conditioned identity admission & 64.12 & 66.66 & +2.54 \\
Admitted source response fusion & 50.59 & 70.38 & +19.79 \\
Native local correspondence & 30.74 & 34.88 & +4.14 \\
Complete-carrier membership reuse & 34.88 & 39.90 & +5.02 \\
\bottomrule
\end{tabular}
}
\caption{Historical matched component controls in our method. The first two
rows measure LERF selected projection and semantic IoU, respectively; the last
two measure UQIS actual image-instance IoU. These interventions are non-additive
and are not final-system leave-one-module-out experiments.}
\label{tab:core-components}
\end{table}
Query-relative identity admission improves projected selection by 2.54 points
with geometry, native features and the set solver fixed. Using admitted source
responses improves semantic IoU by 19.79 points with the evidence, map and
independent selected set fixed. This isolates whether available evidence is
used in the answer, not the novelty of a weighted mean.
Native local image correspondence adds 4.14 points with the same teachers and
extent head; complete-carrier membership adds a further 5.02 in its paired
identity-fixed contrast. The two controls distinguish matching the requested
instance from retrieving all its members.

Additional controls give a 2.35-point category gain from dual-native priors on
three ScanNet rooms, and 9.62-point gains for each UQIS point modality from
real-observation membership training. These include teacher capability and
training-observation changes, respectively. They are not evidence for a
teacher-independent weighting formula; details remain in the supplement.

\subsection{Q2: does completing reference scope before transport recover extent?}
\begin{table}[ht]
\centering\small\begin{tabular}{llrrrr}
\toprule
Matched pipeline & Orig. IoU & Compl. IoU & $\Delta$ & Orig. AP & Compl. AP \\
\midrule
Two-frame transport (N12) & 88.20 & 91.48 & 3.28 & 98.35 & 98.78 \\
Source-only fusion (N15) & 88.01 & 92.75 & 4.74 & 98.26 & 98.81 \\
\bottomrule
\end{tabular}

\caption{Reference-scope interventions on all eight NVOS tasks, retained GT.
Each paired row keeps its downstream propagation, fusion, geometry and readout
fixed. Completion introduces signed reference evidence and DINOv2 into response
construction; it is a mechanism bundle, not an isolated teacher ablation.
Source-only fitting excludes target-image evidence.}
\label{tab:reference-scope}
\end{table}
Under the same frozen two-frame memory and native-ray fitting, replacing the
original reference with the constrained completed reference raises IoU from
88.20 to 91.48 (N12). Semantic completion without the signed constraint yields
88.52; unrestricted semantic similarity does not reliably preserve instance
identity. Reference-only feasibility and quality admission subsequently reach
92.94 and 93.17 under their fixed policies (N14). Their candidates and selection
criteria are independent of target GT; design nevertheless follows feedback
on this opened cohort.

More directly, N15 propagates original and completed references through the
same 318 source views, excludes target-image evidence and applies the same
fusion. IoU rises from 88.01 to 92.75, with AP rising from 98.26 to 98.81.
Five tasks improve and three regress; completion is not uniformly beneficial.
This paired result shows that the scope intervention survives source-view
fusion and is not explained solely by decoding the target image. It does not
separately identify the contribution of negative scribbles, DINOv2 and frozen
SAM3 completion. The central supported claim is about the completed-reference
mechanism under this declared access contract.

\subsection{Q3: are stronger fitting or additional decoding sufficient?}
With fixed evidence, squared reconstruction reduces its source objective on
all 22 source-assisted jobs. Yet LERF semantic IoU/AP change from 70.40/87.62
to 70.80/87.00, and ScanNet mIoU19 falls from 61.01 to 60.28. UQIS outputs are
unchanged compatibility references. Better fit to admitted labels therefore
does not establish better task responses across outputs.

Within the completed-reference experiment, adding target evidence to the source
fit gives 93.03 IoU and 98.98 AP. Against the target-only N14 response,
source residual decreases from 0.004626 to 0.000748 while IoU drops by 0.14
points and AP gains 0.03. All tasks reduce both true and false positives:
consistency suppresses errors but also useful extent. This is a tradeoff,
not evidence that a lower residual solves scope correctness.

Query-local target re-decoding gives 93.28 IoU, only 0.1075 above its N14
control, while AP loses 0.4787; four tasks improve and four regress.
Removing the horns scene cluster eliminates the mean IoU advantage.
A numerical audit exactly recovers all 24 image arrays and checks stable
probability-mixture ranking: the largest macro AP correction is below 0.002
points and no hard masks change. Thus the image AP loss is not explained by
sigmoid saturation or interpolation precision. This audit does not rerun the
3D fits. The joint-fusion and re-decoding search routes are stopped rather than
promoted on one metric.

Finally, retained field-by-readout controls report original NVOS semantic IoU
86.51 and independent selected-set projection IoU 85.22. Crossed outputs are
scored separately. The analysis in Section~\ref{sec:analysis} explains why one
cannot replace the semantic result with a selected-set number or infer better
actual selection from a modified score field. The supplements retain all
crossed readouts, solver failures, prompt controls and per-task diagnostics.

\section{Discussion and limitations}
\paragraph{What the mechanism explains.}
The supported interpretation is that response quality depends on the
applicability and scope of its evidence, as well as its geometric integration.
Identity admission, fixed-evidence response-use controls and reference
completion test different links in that process. The source-only reference
intervention is especially informative: it improves extent before any
target-image fit. Lower reconstruction loss alone and further target decoding
do not consistently improve the resulting answers. These findings support
estimating evidence authority before aggregation, while preserving complete
membership and the intended readout.

\paragraph{What remains an assumption.}
Support, exclusion and candidate quality are estimated. The ambiguity result
does not prove that our completion resolves every unknown region; the exact
attenuation formula is limited to direct consensus with fixed weights; the
occlusion counterexample establishes output non-equivalence, not superiority
of one output for every task. We neither introduce a new weighted-average
primitive nor guarantee monotonic accuracy from admitting evidence. Valid
negative evidence is necessary: removing signed dense hints loses 2.63 IoU
points against its matched NVOS control. Treating all exclusions as invalid
would contradict both the task and the measurements.

\paragraph{Scope of the empirical contribution.}
The five outputs use typed realizations and different legal observation
contracts. Reference completion is measured on NVOS; it has not been shown to
transfer as an unchanged algorithm to all other query types. Native teachers,
source RGB, provided mesh geometry and independent membership supervision are
disclosed, so external score differences are not isolated mechanism effects.
The benchmark cohort has been repeatedly opened during development.
Historical controls are useful causal evidence for their particular contrasts,
not a full factorial ablation of the final system or an independent test of
generalization. The optional learned response selector remains separate.

\paragraph{Remaining validation priorities.}
The next evidence should test the frozen authority/scope rule on unseen scenes
and teacher configurations, and isolate the signed-prompt, completion and
reference-admission contributions under matched access. Failure analysis should
report incomplete extent and wrong-instance expansion separately. More
threshold, crop or fitting sweeps on the current NVOS cohort are a lower
priority: the measured routes have reached their stopping criteria, not a
proven task-wide performance ceiling.

\section{Conclusion}
We study the formation of complete 3D query responses from partial native
teacher evidence. Query-conditioned evidence fusion separates identity
applicability from spatial support and exclusion, then uses geometric or
membership operators matched to the output. Scoped analysis explains why
missing scope, false exclusion and readout mismatch cannot be repaired by
stronger source fitting alone. Five-output evaluation and matched reference
interventions support this mechanism within the declared input and supervision
contracts, while leaving universal cross-teacher transfer open.

\begin{thebibliography}{9}
\bibitem{vala}
S. Wang, K. Li, S. Liang, E. Alegret, J. Ma, N. Navab, and S. Gasperini.
\emph{Visibility-Aware Language Aggregation for Open-Vocabulary Segmentation
in 3D Gaussian Splatting}. arXiv:2509.05515v1, 2025.
\url{https://arxiv.org/abs/2509.05515v1}.
Comparison values are transcribed from Tables 1 and 2; their cited
baselines retain the reporting paper's provenance.
\bibitem{ludvig}
J. Marrie, R. Menegaux, M. Arbel, D. Larlus, and J. Mairal.
\emph{LUDVIG: Learning-Free Uplifting of 2D Visual Features to Gaussian
Splatting Scenes}. ICCV, 2025; arXiv:2410.14462v5.
\url{https://arxiv.org/abs/2410.14462v5}.
NVOS comparison values are transcribed from Tables 1 and D.
\bibitem{occam}
J. Cheng, J.-N. Zaech, L. Van Gool, and D. P. Paudel.
\emph{Occam's LGS: An Efficient Approach for Language Gaussian Splatting}.
arXiv:2412.01807v2, 2025.
\url{https://arxiv.org/abs/2412.01807v2}.
\bibitem{evisplat}
S. Moon, K. Shimomura, J. Park, W. Choi, S. Lee, T. Yamashita, and S. Im.
\emph{EviSplat: Preserving Multi-View Evidence in 3D Gaussian Splatting for
Open-Vocabulary Segmentation}. arXiv:2609.34853v1, 2026.
\url{https://arxiv.org/abs/2609.34853v1}.
\bibitem{sam3}
N. Carion et al. \emph{SAM 3: Segment Anything with Concepts}.
arXiv:2511.16719, 2025. \url{https://arxiv.org/abs/2511.16719}.
\bibitem{siglip2}
M. Tschannen et al. \emph{SigLIP 2: Multilingual Vision-Language Encoders with
Improved Semantic Understanding, Localization, and Dense Features}.
arXiv:2502.14786, 2025. \url{https://arxiv.org/abs/2502.14786}.
\bibitem{dinov2}
M. Oquab et al. \emph{DINOv2: Learning Robust Visual Features without Supervision}.
arXiv:2304.07193, 2023. \url{https://arxiv.org/abs/2304.07193}.
\bibitem{maskadapter}
Y. Li, T. Cheng, B. Feng, W. Liu, and X. Wang.
\emph{Mask-Adapter: The Devil is in the Masks for Open-Vocabulary Segmentation}.
CVPR, 2025. \url{https://github.com/hustvl/MaskAdapter}.
\bibitem{vcar}
K. Cao, D. Wang, H. Zhu, H. Huang, X. Wang, Z. Shi, and G. Yang.
\emph{VCAR: Training-Free 3DGS Segmentation via View Completeness and Axis-Aware Boundary Refinement}.
arXiv:2608.30870v1, 2026. \url{https://arxiv.org/abs/2608.30870v1}.
\bibitem{lego}
Y. Peng, H. Wang, Y. Liu, Y. Lu, Z. Dong, and B. Yang.
\emph{LEGO: Leveled Language Gaussian Splatting}.
ECCV, 2026. \url{https://arxiv.org/abs/2608.10057v1}.
\bibitem{wildseg3d}
Y. Guo, J. Hu, Y. Qu, and L. Cao.
\emph{WildSeg3D: Segment Any 3D Objects in the Wild from 2D Images}.
ICCV, 2025. \url{https://arxiv.org/abs/2503.08407}.
\bibitem{trace3d}
H. Shen, J. Ni, Y. Chen, W. Li, M. Pei, and S. Huang.
\emph{Trace3D: Consistent Segmentation Lifting via Gaussian Instance Tracing}.
ICCV, 2025. \url{https://arxiv.org/abs/2508.03227v1}.
\end{thebibliography}

\end{document}
